\begin{exo}
	Dans une boulangerie, on dispose d'une réduction si l'on choisit la formule \og{} dessert mystère \fg{} pour laquelle le dessert accompagnant le menu est tiré au hasard.

	Ceyda choisit cette formule alors que les desserts encore disponibles sont répartis comme suit :

	\begin{center}
		\renewcommand{\arraystretch}{1.2}
		\begin{tabular}{|c|c|c|c|} \hline
			           & Chocolat     & Vanille     & Total        \\ \hline
			Tartelette & \phantom{8}  & 11          & 19           \\ \hline
			Éclair     & \phantom{13} & \phantom{7} & \phantom{20} \\ \hline
			Total      & \phantom{21} & 18          & 39           \\ \hline
		\end{tabular}
	\end{center}

	On considère les événements $E$ : \og{} Son dessert est un éclair \fg{} et $V$ : \og{} Son dessert est à la vanille. \fg{}

	\begin{enumerate}
		\item Calculer $p_{E}(V)$, $p_{V}(E)$, $p_{\overline{E}}(V)$.
		\item Ceyda voit que son dessert est un éclair.

		      Écrire la probabilité qu'il soit au chocolat comme une probabilité conditionnelle puis la calculer.
	\end{enumerate}
\end{exo}
